\documentclass[11pt, a4paper]{article}

\usepackage[T1]{fontenc}
\usepackage[utf8]{inputenc}
\usepackage{tgpagella}
\usepackage[spanish, es-noshorthands]{babel}
\usepackage{amsmath, amssymb, bm}
\usepackage{tabularx}
\usepackage{booktabs}
\usepackage[most]{tcolorbox}
\usepackage[margin=2.5cm]{geometry}
\usepackage{ragged2e}
\usepackage{fancyhdr}
\usepackage[american]{circuitikz}

\pagestyle{fancy}
\fancyhf{}
\setlength{\headheight}{16pt}
\lhead{\textbf{UCB Sede Tarija} | Ing. Mecatrónica}
\chead{}
\rhead{IMT-121 Circuitos Electrónicos I | LTspice RLC}
\lfoot{Prof: M.Sc. Ing. Bernardo Quiroga Turdera}
\rfoot{Página \thepage}
\renewcommand{\headrulewidth}{0.4pt}

\definecolor{ucbblue}{RGB}{0, 51, 102}

\begin{document}

\begin{tcolorbox}[colback=ucbblue!5!white, colframe=ucbblue, title=\textbf{Práctica de Simulación y Análisis de Discrepancias: Red RLC Serie}]
    \centering \textbf{Uso de Barridos Paramétricos (.step param)}
\end{tcolorbox}

\textbf{Estudiante:} \hrulefill \quad \textbf{Fecha:} \hrulefill

\vspace{0.3cm}
\begin{tcolorbox}[colback=gray!5, colframe=gray!50, title=\textbf{Modalidad de Evidencia Tecnológica}]
\textit{Marque su escenario de ejecución para mantener trazabilidad:} \par\vspace{0.2cm}
{[$\square$]} Ejecución Indiv. \quad {[$\square$]} Ejecución Compartida \quad {[$\square$]} Demostración \quad {[$\square$]} Acceso Pendiente.
\end{tcolorbox}

\section*{1. Representación del Circuito en SPICE}
Dibuje el siguiente modelo de simulación utilizando $L=10\,\text{mH}$ y $C=1\,\mu\text{F}$. Para evaluar los 3 regímenes simultáneamente, parametrizaremos la resistencia total.
\begin{itemize}
    \item Parámetros fuente: \texttt{PULSE(0 10 100u 1u 1u 5m 10m)}
    \item Resistencia parametrizada: Valor $\to$ \texttt{\{Rtot\}}
    \item Directivas a incrustar en el plano: \\
          \texttt{.param Rtot=50} \\
          \texttt{.step param Rtot list 50 200 400} \\
          \texttt{.tran 0 3m 0 2u}
\end{itemize}
\textit{Nota: No use comandos \texttt{uic} ni inicialice condensadores artificialmente. SPICE calculará el punto de operación pasivo en $0\,\text{V}$ antes del flanco de los $100\,\mu\text{s}$.}

\section*{2. Tabla de Contraste Experimental}
Ejecute la simulación. Añada la traza del voltaje en el nodo del capacitor \texttt{V(out)} y complétela:
\vspace{0.3cm}

\renewcommand{\arraystretch}{1.5}
\noindent
% Al usar \hsize matemáticamente distribuido, la suma es 6 (0.8*4 + 1.2 + 1.6 = 6). Garantiza ajuste exacto al \textwidth sin desbordes.
\begin{tabularx}{\textwidth}{
  >{\centering\arraybackslash\hsize=0.8\hsize}X 
  >{\centering\arraybackslash\hsize=0.8\hsize}X 
  >{\RaggedRight\arraybackslash\hsize=1.2\hsize}X 
  >{\centering\arraybackslash\hsize=0.8\hsize}X 
  >{\centering\arraybackslash\hsize=0.8\hsize}X 
  >{\RaggedRight\arraybackslash\hsize=1.6\hsize}X}
\toprule
\textbf{$R_{total}$ Nom.} & \textbf{$\zeta$ Teórico} & \textbf{Régimen Predicho} & \textbf{¿Hay Sobreimpulso?} & \textbf{Lectura Pico Máximo} & \textbf{Observación Causal} \\ 
\midrule
$50\,\Omega$ & $0.25$ & & & & \\ \addlinespace
$200\,\Omega$ & $1.00$ & & & & \\ \addlinespace
$400\,\Omega$ & $2.00$ & & & & \\
\bottomrule
\end{tabularx}

\section*{3. Discusión de Simulación vs. Mundo Real}
Responda brevemente:
\begin{enumerate}
    \item ¿Por qué $v_C(0^+) = 0\,\text{V}$ incluso en la curva azul (Subamortiguada) que supera brutalmente los $10\,\text{V}$? \par
    \vspace{0.2cm}
    \noindent\hrulefill
    
    \item ¿Qué ocurriría si intentamos verificar esto en protoboard utilizando un resistor exacto de $200\,\Omega$ asumiendo que es crítico? ¿Qué componentes "ocultos" alterarán $\zeta$? \par
    \vspace{0.2cm}
    \noindent\hrulefill
\end{enumerate}

\end{document}
